\documentclass[12pt,a4paper,reqno]{amsart}

% language
\usepackage[greek,english]{babel}
\usepackage[utf8]{inputenc}
\usepackage{alphabeta}

% change default names to greek
\addto\captionsenglish{
    \renewcommand{\contentsname}{Περιεχόμενα}
    \renewcommand{\refname}{Βιβλιογραφία}
    \renewcommand{\datename}{Ημερομηνία:}
    \renewcommand{\urladdrname}{Ιστοσελίδα}
}

% math 
\usepackage{amsmath,amsthm,amssymb,amscd}

% font
\usepackage[cal=euler]{mathalfa}
\usepackage{libertinus-type1}
% \usepackage{txfonts} % for upright greek letters
\usepackage{bm} % for bold symbols
\usepackage{bbm} % for the simply-looking bb symbols

% miscellaneous 
\usepackage{changepage} %for indenting environments
\usepackage{csquotes} % example: \textcquote{}
\usepackage{blkarray}
\setcounter{MaxMatrixCols}{20} % default for pmatrix is 10!!
\usepackage{ytableau}
\usepackage{array} %needed to increase the vertical length in a tabular

% drawing
\usepackage{tikz,tikz-cd}
\usetikzlibrary{shapes.misc, patterns, matrix, calc, intersections,positioning,arrows.meta}
\usepackage{graphics,graphicx}
\usepackage{float} % provides enhanced control and customization options for floating objects such as figures and tables

% colors
\usepackage{xcolor}
\definecolor{darkcandyapplered}{rgb}{0.64, 0.0, 0.0}
\definecolor{midnightblue}{rgb}{0.1, 0.1, 0.44}
\definecolor{mylightblue}{HTML}{336699}
\definecolor{burntorange}{rgb}{0.8, 0.33, 0.0}
\definecolor{iceberg}{rgb}{0.44, 0.65, 0.82}
\definecolor{applegreen}{rgb}{0.55, 0.71, 0.0}
\definecolor{canaryyellow}{rgb}{1.0, 0.94, 0.0}
\definecolor{lightgray}{rgb}{0.83, 0.83, 0.83}

% hrefs
\usepackage{hyperref}
\usepackage[noabbrev,capitalize]{cleveref}
\hypersetup{
    pdftoolbar=true,        
    pdfmenubar=true,        
    pdffitwindow=false,     
    pdfstartview={FitH},  % fits the width of the page to the window
    pdftitle={},
    pdfauthor={},
    pdfsubject={},
    pdfkeywords={},
    pdfnewwindow=true,  % links in new window
    colorlinks=true,  % false: boxed links; true: colored links
    linkcolor=darkcandyapplered,   % color of internal links
    citecolor=midnightblue,  % color of links to bibliography
    urlcolor=cyan,  % color of external links
    linktocpage=true  % changes the links from the section body to the page number
    }

% geometry
\textwidth=16cm 
\textheight=21cm 
\hoffset=-55pt 
\footskip=25pt

% thm envs (you might need to change the path)
\theoremstyle{definition}
\newtheorem{exercise}{Άσκηση}
\newtheorem{solution}{Λύση}
\newtheorem*{remark}{Παρατήρηση}
\newtheorem*{example}{Παράδειγμα}

% math ops 
% fields
\newcommand{\NN}{\mathbbmss{N}} 
\newcommand{\ZZ}{\mathbbmss{Z}} 
\newcommand{\QQ}{\mathbbmss{Q}} 
\newcommand{\RR}{\mathbbmss{R}} 
\newcommand{\CC}{\mathbbmss{C}} 
\newcommand{\KK}{\mathbbmss{K}} 
\newcommand{\FF}{\mathbbmss{F}} 
\newcommand{\HH}{\mathbbmss{H}} 

% symmetric group
\newcommand{\fS}{\mathfrak{S}}  
\newcommand{\fp}{\mathfrak{p}}  
\newcommand{\fm}{\mathfrak{m}}  

% calligraphic 
\newcommand{\aA}{\mathcal{A}} 
\newcommand{\bB}{\mathcal{B}}
\newcommand{\cC}{\mathcal{C}}
\newcommand{\dD}{\mathcal{D}}
\newcommand{\eE}{\mathcal{E}}
\newcommand{\fF}{\mathcal{F}}
\newcommand{\hH}{\mathcal{H}}
\newcommand{\iI}{\mathcal{I}}
\newcommand{\lL}{\mathcal{L}}
\newcommand{\nN}{\mathcal{N}}
\newcommand{\oO}{\mathcal{O}}
\newcommand{\pP}{\mathcal{P}}
\newcommand{\sS}{\mathcal{S}}
\newcommand{\mM}{\mathcal{M}}
\newcommand{\uU}{\mathcal{U}}

% bold
\newcommand{\bfa}{\mathbf{a}}
\newcommand{\bfe}{\mathbf{e}}
\newcommand{\bfF}{\pmb{F}}
\newcommand{\bfR}{\pmb{R}}
\newcommand{\bfv}{\mathbf{v}}
%\newcommand{\bfx}{\bm{x}}
%\newcommand{\bfx}{\mathbf{x}} 
\newcommand{\bfx}{\pmb{x}}
\newcommand{\bfX}{\pmb{X}}
\newcommand{\bfy}{\pmb{y}}
\newcommand{\bfz}{\pmb{z}}

% roman
\newcommand{\rmA}{\mathrm{A}}
\newcommand{\rmB}{\mathrm{B}}
\newcommand{\rmC}{\mathrm{C}}
\newcommand{\rmD}{\mathrm{D}} 
\newcommand{\rmH}{\mathrm{H}} 
\newcommand{\rmh}{\mathrm{h}} 
\newcommand{\rmI}{\mathrm{I}} 
\newcommand{\rmJ}{\mathrm{J}} 
\newcommand{\rmK}{\mathrm{K}}
\newcommand{\rmM}{\mathrm{M}}
\newcommand{\rmP}{\mathrm{P}}  
\newcommand{\rmp}{\mathrm{p}}  
\newcommand{\rmQ}{\mathrm{Q}}  
\newcommand{\rmR}{\mathrm{R}}
\newcommand{\rmS}{\mathrm{S}}
\newcommand{\rmT}{\mathrm{T}}
\newcommand{\rmU}{\mathrm{U}}
\newcommand{\rmV}{\mathrm{V}}
\newcommand{\rmY}{\mathrm{Y}}
\newcommand{\rmZ}{\mathrm{Z}}
\newcommand{\rmz}{\mathrm{z}}

% arrows and symbols 
\renewcommand{\to}{\rightarrow}
\newcommand{\toto}{\longrightarrow}
\newcommand{\mapstoto}{\longmapsto}
\newcommand{\then}{\Rightarrow}
\newcommand{\IFF}{\Leftrightarrow}
\newcommand{\tl}{\tilde}
\newcommand{\wtl}{\widetilde}
\newcommand{\ol}{\overline}
\newcommand{\ul}{\underline}
\newcommand{\oldemptyset}{\emptyset}
\renewcommand{\emptyset}{\varnothing}
\DeclareMathSymbol{\Arg}{\mathbin}{AMSa}{"39} % for arguments 
\newcommand{\onto}{\ensuremath{\twoheadrightarrow}}
\newcommand{\tle}{\trianglelefteq}
\newcommand{\tge}{\trianglerighteq}

% custom ops
\newcommand{\rmKer}{\mathrm{Ker}}
\newcommand{\rmIm}{\mathrm{Im}}
\newcommand{\lcm}{\mathrm{lcm}}
\newcommand{\GL}{\mathrm{GL}}
\newcommand{\ev}{\mathrm{ev}}
\newcommand{\End}{\mathrm{End}}
\newcommand{\rmid}{\mathrm{id}}
\newcommand{\rmspan}{\mathrm{span}}
\newcommand{\Spec}{\mathrm{Spec}}
\newcommand{\mSpec}{\mathrm{mSpec}}
\newcommand{\rad}{\mathrm{rad}}
\newcommand{\Rad}{\mathrm{Rad}}
\newcommand{\Jac}{\mathrm{Jac}}
\newcommand{\tor}{\mathrm{tor}}
\newcommand{\Ann}{\mathrm{Ann}}
\newcommand{\Hom}{\mathrm{Hom}}

% absolute value symbol
\usepackage{mathtools} 
\DeclarePairedDelimiter\abs{\lvert}{\rvert}%
\DeclarePairedDelimiter\norm{\lVert}{\rVert}%
\makeatletter
\let\oldabs\abs
\def\abs{\@ifstar{\oldabs}{\oldabs*}}

% tensor symbol
\newcommand{\tensor}[1]{%
  \mathbin{\mathop{\otimes}\limits_{#1}}%
}

% permutation cycle notation
\ExplSyntaxOn
\NewDocumentCommand{\cycle}{ O{\;} m }
 {
  (
  \alec_cycle:nn { #1 } { #2 }
  )
 }

\seq_new:N \l_alec_cycle_seq
\cs_new_protected:Npn \alec_cycle:nn #1 #2
 {
  \seq_set_split:Nnn \l_alec_cycle_seq { , } { #2 }
  \seq_use:Nn \l_alec_cycle_seq { #1 }
 }
\ExplSyntaxOff

% setminus symbol
\newcommand{\mysetminusD}{\hbox{\tikz{\draw[line width=0.6pt,line cap=round] (3pt,0) -- (0,6pt);}}}
\newcommand{\mysetminusT}{\mysetminusD}
\newcommand{\mysetminusS}{\hbox{\tikz{\draw[line width=0.45pt,line cap=round] (2pt,0) -- (0,4pt);}}}
\newcommand{\mysetminusSS}{\hbox{\tikz{\draw[line width=0.4pt,line cap=round] (1.5pt,0) -- (0,3pt);}}}
\newcommand{\sm}{\mathbin{\mathchoice{\mysetminusD}{\mysetminusT}{\mysetminusS}{\mysetminusSS}}}

% miscellaneous commands
\newcommand{\defn}[1]{{\color{mylightblue}{#1}}}
\newcommand{\toDo}{{\bf\color{red} TODO}}
\newcommand{\toCite}{{\bf\color{green} CITE}}
\newcommand*{\vertbar}{\rule[-1ex]{0.5pt}{2.5ex}} % for matrices with column vectors
\newcommand*{\horzbar}{\rule[.5ex]{2.5ex}{0.5pt}} % for matrices with row vectors
\newcommand{\myblue}[1]{{\color{iceberg}{#1}}}
\newcommand{\myorange}[1]{{\color{burntorange}{#1}}}
\newcommand{\mygreen}[1]{{\color{applegreen}{#1}}}
\newcommand{\myred}[1]{{\color{darkcandyapplered}{#1}}}

% ferrer's diagram
\newcommand{\fdiagram}[1]{
    \begin{tikzpicture}[scale=.7]
        \fill foreach \Z [count=\Y] in {#1}
        {foreach \X in {1,...,\Z} 
        {(\X,-\Y) circle[radius=3pt]}};
    \end{tikzpicture}
}

%
\usepackage{pgfplots}
\pgfplotsset{compat=1.18}

%
\makeatletter
\def\mathcolor#1#{\@mathcolor{#1}}
\def\@mathcolor#1#2#3{%
  \protect\leavevmode
  \begingroup
    \color#1{#2}#3%
  \endgroup
}
\makeatother

% 
\newenvironment{nouppercase}{%
  \let\uppercase\relax%
  \renewcommand{\uppercasenonmath}[1]{}}{}

% titlepage
\title{226: Θεωρία Δακτυλίων και Modules \\ Ασκήσεις με Ενδεικτικές Λύσεις}
\author[Β.~Δ. Μουστακας]{Βασίλης Διονύσης Μουστάκας \\ Πανεπιστήμιο Κρήτης}
\date{27 Μαρτίου 2026}
% \urladdr{\href{https://sites.google.com/view/vasmous}{https://sites.google.com/view/vasmous}}

\begin{document}

\begingroup
\def\uppercasenonmath#1{} % this disables uppercase title
\let\MakeUppercase\relax % this disables uppercase authors
\maketitle
\endgroup

\setcounter{exercise}{19}
% \setcounter{theorem}{5}
\setcounter{equation}{3}
\renewcommand{\thesubsection}{\thesection.\arabic{subsection}}

Σε ότι ακολουθεί, υποθέτουμε ότι $R$ είναι ένας δακτύλιος και $n$ είναι ένας θετικός ακέραιος, εκτός αν αναφέρεται διαφορετικά.

\begin{exercise}{\rm(\cite[Ασκήσεις~10.2.9~και~10.2.10]{DF04})}
  \label{ex:DF_10.2.9_10.2.10}
  Υποθέτουμε ότι ο $R$ είναι μεταθετικός δακτύλιος και έστω $M$ ένα $R$-πρότυπο. Να δείξετε ότι 
  \begin{itemize}
    \item[(1)] $\Hom_R(R,M) \cong M$ ως $R$-πρότυπα, και 
    \item[(2)] $\End_R(R) \cong R$ ως δακτύλιοι.
  \end{itemize}
\end{exercise}

\begin{proof}[Λύση]
  Για το (1), θεωρούμε τις απεικονίσεις
  \begin{align*}
    \Phi : \Hom_R(R,M) &\to M \\
    f &\mapsto f(1_R)
  \end{align*}
  και
  \begin{align*}
    \Psi : M &\to \Hom_R(R,M) \\
            m &\mapsto f_m,
  \end{align*}
  όπου $f_m(r) = rm$. Αυτές είναι $R$-γραμμικές και ικανοποιούν 
  \[
  \Phi \circ \Psi = \rmid_M \quad \text{και} \quad 
  \Psi \circ \Phi = \rmid_{\Hom_R(M,N)}.
  \]
  Για το (2), αρκεί να παρατηρήσουμε ότι οι απεικονίσεις $\Phi$ και $\Psi$ είναι ομομορφισμοί δακτυλίων. Οι λεπτομέρειες αφήνονται ως άσκηση.
\end{proof}

\begin{exercise}{\rm(\cite[Άσκηση~10.2.11]{DF04})}
  \label{ex:DF_10.2.11}
  Για κάθε $1 \le i \le n$, Θεωρούμε ένα $R$-πρότυπο $A_i$ και ένα υποπρότυπο αυτού $B_i$. Να δείξετε ότι 
  \[
  \frac{A_1\times A_2 \times \cdots \times A_n}{B_1 \times B_2 \times \cdots \times B_n} \cong \frac{A_1}{B_1} \times \frac{A_2}{B_2} \times \cdots \times \frac{A_n}{B_n}.
  \]
\end{exercise}

\begin{proof}[Λύση]
  Θεωρούμε την απεικόνιση 
  \begin{align*}
    \phi : A_1 \times A_2 \times \cdots \times A_n &\to \frac{A_1}{B_1} \times \frac{A_2}{B_2} \times \cdots \times \frac{A_n}{B_n} \\
    (a_1, a_2, \dots, a_n) &\mapsto (a_1 + B_1, a_2 + B_2, \dots, a_n + B_n).
  \end{align*}
  Η $\phi$ είναι $R$-επιμορφισμός και 
  \[
  \ker(\phi) = B_1 \times B_2 \times \cdots \times B_n
  \]
  (γιατί;). Το ζητούμενο έπεται από το πρώτο θεώρημα ισομορφισμού.
\end{proof}

\begin{exercise}{\rm(\cite[Άσκηση~10.3.7]{DF04})}
  \label{ex:DF_10.3.7}
  Έστω
  \begin{equation}
    \label{eq:DF_10.3.7}
    \begin{tikzcd}
        0 \arrow[r] & X \arrow[r, "\alpha"] & Y \arrow[r, "\beta"] & Z \arrow[r] & 0.
    \end{tikzcd}
  \end{equation}
  μία βραχεία ακριβής ακολουθία. Να δείξετε ότι αν τα $X$ και $Z$ είναι πεπερασμένα παραγόμενα, τότε και το $Y$ είναι πεπερασμένα παραγόμενο. Ειδικότερα, συνάγετε ότι κάθε $R$-πρότυπο $M$ για το οποίο κάθε υποπρότυπο $N$ και κάθε $M/N$ είναι πεπερασμένα παραγόμενα είναι πεπερασμένα παραγόμενο.
\end{exercise}

\begin{proof}[Λύση]
  Ο δεύτερος ισχυρισμός έπεται εφαρμόζοντας τον πρώτο στην βραχεία ακριβή ακολουθία 
  \[
    \begin{tikzcd}
        0 \arrow[r] & N \arrow[r, hook] & \underbrace{M \oplus (M/N)}_{\cong \ M} \arrow[r, two heads] & M/N \arrow[r] & 0
    \end{tikzcd}
  \]
  (γιατί;).
  
  Για τον πρώτο ισχυρισμό, έστω $X = R\{x_1, x_2, \dots, x_n\}$ και $Z = R\{z_1, z_2, \dots, z_m\}$. Θα δείξουμε ότι 
  \[
  Y = R\{\alpha(x_1), \alpha(x_2), \dots, \alpha(x_n), z_1, z_2, \dots, z_m\}.
  \] 
  Επειδή η ακολουθία \eqref{eq:DF_10.3.7} είναι ακριβής στο $Z$ έπεται ότι η $\beta$ είναι $R$-επιμορφισμός και γι αυτό μπορούμε να διαλέξουμε $y_i \in Y$ τέτοια ώστε 
  \[
  \beta(y_i) = z_i,
  \]
  για κάθε $1 \le i \le m$. 

  Από την άλλη, για κάθε $y \in Y$ έχουμε 
  \begin{align*}
  \beta(y) &= r_1z_1 + r_2z_2 + \cdots + r_mz_m \\ 
  &= r_1\beta(y_1) + r_2\beta(y_2) + \cdots + r_m\beta(y_m) \\
  &= \beta(r_1y_1) + \beta(r_2y_2) + \cdots + \beta(r_my_m) \\
  &= \beta\left(
    r_1y_1 + r_2y_2 + \cdots + r_my_m
  \right),
  \end{align*}
  για κάποια $r_1, r_2, \dots, r_m \in R$. Με άλλα λόγια, 
  \[
  \beta\left(
    y - \left(
    r_1y_1 + r_2y_2 + \cdots + r_my_m
    \right)
  \right) = 0_Z
  \]
  ή ισοδύναμα 
  \[
  y - \left(
    r_1y_1 + r_2y_2 + \cdots + r_my_m
    \right) \in \ker(\beta).
  \]

  Επειδή η ακολουθία \eqref{eq:DF_10.3.7} είναι ακριβής στο $Y$, έπεται ότι 
  \[
  y - \left(
    r_1y_1 + r_2y_2 + \cdots + r_my_m
    \right) = s_1\alpha(x_1) + s_2\alpha(x_2) + \cdots + s_n\alpha(x_n)
  \]
  για κάποια $s_1, s_2, \dots, s_n \in R$ (γιατί;) και γι αυτό
  \[
  y = r_1y_1 + r_2y_2 + \cdots + r_my_m + s_1\alpha(x_1) + s_2\alpha(x_2) + \cdots + s_n\alpha(x_n)
  \]
  και το ζητούμενο έπεται.
\end{proof}

\begin{remark}
  Το αντίστροφο του δεύτερου ισχυρισμού δεν ισχύει πάντα. Για παράδειγμα, έστω 
  \[
  R = M = \ZZ[x_1, x_2, \dots]
  \]
  ο πολυωνυμικός δακτύλιος σε (αριθμήσιμα) άπειρες μεταβλητές $x_1, x_2, \dots$. Αυτός είναι πεπερασμένα παραγόμενος (όπως και κάθε κανονικό πρότυπο), αλλά το υποπρότυπο 
  \[ 
  N = (x_1, x_2, \dots)
  \] 
  δεν είναι (γιατί;).
\end{remark}

\begin{exercise}{\rm(\cite[Άσκηση~10.3.9]{DF04})}
  \label{ex:DF_10.3.9}
  Ένα μη μηδενικό $R$-πρότυπο $M$ ονομάζεται \defn{απλό} (irreducible ή simple) αν τα μόνα υποπρότυπά του είναι το $0$ και το $M$.
  \begin{itemize}
    \item[(1)] Να δείξετε ότι το $M$ είναι απλό αν και μόνο αν είναι κυκλικό και παράγεται από κάθε μη μηδενικό στοιχείο.
    \item[(2)] Να καθορίσεται όλα τα απλά $\ZZ$-πρότυπα.
  \end{itemize}
\end{exercise}

\begin{proof}[Λύση]
  \leavevmode
  \begin{itemize}
    \item[(1)] Υποθέτουμε ότι το $M$ είναι απλό $R$-πρότυπο και θεωρούμε ένα μη μηδενικό $m \in M$. Τότε, το υποπρότυπο $Rm$ είναι μη μηδενικό και κατά συνέπεια $Rm = M$.
    
    Αντιστρόφως, υποθέτουμε ότι το $M$ παράγεται από κάθε μη μηδενικό στοιχείο και θεωρούμε ένα υποπρότυπο $N$ του $M$. Αν $N \neq 0$, τότε υπάρχει μη μηδενικό $a \in N \subseteq$ για το οποίο $M = Ra \subseteq N$ και κατά συνέπεια $M = N$. Διαφορετικά, $N=0$.
    \item[(2)] Από το Ερώτημα (1), τα πιθανά απλά $\ZZ$-πρότυπα είναι οι κυκλικές αβελιανές ομάδες, δηλαδή $\ZZ$ και $\ZZ_n$. 
    \begin{itemize}
      \item Παρατηρούμε ότι 
      \[
      (2) = \{0, 2, 4, \dots\} \neq \ZZ.
      \]
      και γι αυτό, από το Ερώτημα (1), το $\ZZ$ δεν μπορεί να είναι απλό.
      \item Από τη θεωρία ομάδων, γνωρίζουμε ότι ένα στοιχείο $k \in \ZZ_n$ παράγει την $\ZZ_n$ αν και μόνο άν $\gcd(k,n) = 1$. Από το Ερώτημα (1), αυτό συμβαίνει για κάθε μη μηδενικό στοιχεία αν και μόνο αν το $n$ είναι πρώτος.
    \end{itemize}
    Κατά συνέπεια, τα απλά $\ZZ$-πρότυπα είναι τα $\ZZ_p$ για κάποιο πρώτο $p$.
  \end{itemize}
\end{proof}

\begin{exercise}{\rm(\cite[Άσκηση~10.3.10]{DF04})}
  \label{ex:DF_10.3.10}
  Υποθέτουμε ότι ο $R$ είναι μεταθετικός δακτύλιος.
  \begin{itemize}
    \item[(1)] Να δείξετε ότι το $M$ είναι απλό αν και μόνο αν $M \cong R/I$ για κάποιο $I \in \mSpec(R)$.
    \item[(2)] Να καθορίσεται όλα τα απλά $\ZZ$-πρότυπα (ξανά).
  \end{itemize}
\end{exercise}

\begin{proof}[Λύση]
  \leavevmode
  \begin{itemize}
    \item[(1)] Υποθέτουμε ότι το $M$ είναι απλό. Από την Άσκηση \ref{ex:DF_10.3.9} έπεται ότι το $M$ είναι κυκλικό και παράγεται από κάθε μη μηδενικό στοιχείο. Επειδή ο $R$ είναι μεταθετικός, από το Παράδειγμα 4.14 (δ) έπεται ότι $M \cong R/I$ για κάποιο ιδεώδες\footnote{Για την ακρίβεια, αν $M = (m)$, για κάποιο $m \neq 0$, τότε 
    $I = \Ann_R((m)) = \{r \in R : rm = 0_M \}$.} $I$. Επομένως, τα υποπρότυπα του $M$ αντιστοιχούν στα ιδεώδη του $R/I$ με τα τελευταία να αντιστοιχούν στα ιδεώδη $J$ του $R$ για τα οποία $I \subseteq J$, από το τέταρτο θεώρημα ισομορφισμού. Επειδή το $M$ είναι απλό αυτό σημαίνει ότι είτε $J = I$ είτε $J = R$ και κατά συνέπεια $J \in \mSpec(R)$.

    Αντιστρόφως, υποθέτουμε ότι $M \cong R/I$ για κάποιο $I \in \mSpec(R)$. Από την Πρόταση 1.10, το $M$ είναι σώμα και κατά συνέπεια απλό $R$-πρότυπο.
    \item[(2)] Γνωρίζουμε ότι $\mSpec(\ZZ) = \{(p) : p \ \text{πρώτος}\}$. Από το Ερώτημα (1) έπεται ότι τα απλά $\ZZ$-πρότυπα είναι της μορφής $\ZZ_p$, για κάποιο πρώτο $p$, σε συμφωνία με τον υπολογισμό της Άσκησης \ref{ex:DF_10.3.9}.
  \end{itemize}
\end{proof}

\begin{remark}
  Αν κάθε $R$-πρότυπο μπορεί να γραφεί ως ευθύ άθροισμα απλών προτύπων, τότε ο $R$ ονομάζεται \defn{ημιαπλός} (semisimple). Έστω $G$ μια πεπερασμένη ομάδα και $F$ ένα σώμα. Αποδεικνύεται ότι ο ομαδοδακτύλιος $FG$ της $G$ (που ορίστηκε στην Άσκηση 2) είναι ημιαπλός όταν το $F$ είναι σώμα χαρακτηριστικής 0 (όπως, για παράδειγμα, το $\CC$). Η μελέτη των $FG$-προτύπων ονομάζεται \emph{θεωρία αναπαραστάσεων} πεπερασμένων ομάδων (βλ., για παράδειγμα, \cite[Μέρος~VI]{DF04}).
\end{remark}

\begin{exercise}{\rm(Λήμμα του Schur \cite[Άσκηση~10.3.11]{DF04})}
  \label{ex:DF_10.3.11}
  Να δείξετε τα εξής.
  \begin{itemize}
    \item[(1)] Αν $M, N$ είναι δύο απλά $R$-πρότυπα και $\phi \in \Hom_R(M,N)$, τότε είτε $\phi = 0$, είτε η $\phi$ είναι ισομορφισμός.
    \item[(2)] Αν $M$ είναι απλό $R$-πρότυπο, τότε ο $\End_R(M)$ είναι δακτύλιος διαίρεσης.
  \end{itemize}
\end{exercise}

\begin{proof}[Λύση]
  Το (2) έπεται άμεσα από το (1), διότι κάθε μη μηδενικός $R$-ενδομορφισμός $\phi : M \to M$ είναι $R$-ισομορφισμός και γι αυτό $\phi \circ \phi^{-1} = \phi^{-1}\circ \phi = \rmid_M$, δηλαδή αντιστρέψιμος. Για το (1), επειδή το $M$ είναι απλό έπεται ότι 
  \[
  \ker(\phi) = 0_M \quad \text{ή} \quad \ker(\phi) = M.
  \]
  Στη δεύτερη περίπτωση, έχουμε αναγκαστικά $\phi = 0$ (γιατί;). Διαφορετικά, επειδή το $N$ είναι απλό έπεται ότι 
  \[
  \rmIm(\phi) = 0_N \quad \text{ή} \quad \rmIm(\phi) = N.
  \]
  Στην πρώτη περίπτωση, έχουμε αναγκαστικά $\phi = 0$ (γιατί;). Διαφορετικά, η $\phi$ είναι $R$-ισομορφισμός.
\end{proof}
 
\begin{thebibliography}{99}
    \bibitem{DF04}
    D. S. Dummit και R. M. Foote, 
    \textit{Abstract algebra},
    third edition, Wiley, 2004. 
    % \bibitem{Mal08}
    % Μ. Μαλιάκας, 
    % \textit{Εισαγωγή στην Μεταθετική Άλγεβρα}, 
    % Εκδόσεις Σοφία, 2008.
\end{thebibliography}
\end{document}